\section{Discussion}\label{sec:discussion}

\subsection{Main inference from the ablation}
\label{subsec:discussion-main-inference}

The ablation results indicate that the principal source of improvement
is not stochastic perturbation alone, nor the nominal complexity of the
ML backend, but the transfer of distributional information from the
stochastic physical ensemble to the emulator. This distinction is
important for interpreting the experiment hierarchy. The deterministic
and stochastic physical baselines remain within a narrow performance
range, and direct L\'{e}vy perturbation of the baseflow-enabled physical
model does not improve point-prediction skill. By contrast, the
backend-consistent quantile hybrid E6 improves NSE by 0.311 relative to
the full stochastic physical baseline E3 and by 0.165 relative to the
pure forcing-based ML baseline E4. The most direct descriptor ablation,
E6--E5, adds 0.114 NSE by replacing an ensemble-mean feature with a
vector of ensemble quantiles while preserving the same executed
backend. This result supports a specific interpretation of hybridization:
the physical model is most useful when it is used to construct
informative state-dependent predictors, rather than when its stochastic
output is treated as a calibrated forecast distribution.

This interpretation is consistent with recent reviews of hydrological
ML and hybrid modelling, which emphasize that performance gains can
come either from more flexible algorithms or from more informative
predictors \citep{papacharalampous2022review,zhong2023developing}. In
StoHyMoLAP, the ablation evidence points mainly to the second mechanism.
The stochastic physical ensemble transforms precipitation,
evapotranspiration and antecedent basin conditions into a compact set of
hydrologically filtered descriptors. The ML model then learns a residual
or corrective mapping from these descriptors to observed discharge. The
result is therefore not evidence that ML replaces hydrological
structure; rather, it shows that a simple physical structure can become
more useful when its distributional response is exposed to the emulator
\citep{mekonnen2015hybrid,houenafa2025hybridization}.

\subsection{Hydrological memory, baseflow and low-flow behaviour}
\label{subsec:discussion-baseflow}

The E1--E0 contrast shows that adding a linear baseflow reservoir
produces a small but positive global gain in the deterministic physical
baseline. This is hydrologically plausible for an Altiplano basin with
strong wet--dry seasonality, delayed drainage and low-flow persistence.
Previous Andean and mountain-catchment studies have similarly shown
that groundwater storage, snow/glacier contributions, wetlands and
season-dependent parameter behaviour can control dry-season discharge
and recession dynamics \citep{nauditt2017conceptual,ragettli2014evaluation,condom2004transient,lima2025modeling}. The baseflow store in
StoHyMoLAP should therefore be interpreted as a minimal representation
of hydrological memory rather than as a complete description of
subsurface storage.

The regime-specific results also show the limitation of this minimal
representation. Although E1 improves the global deterministic score, the
low-flow RMSE in Table~\ref{tab:regime-rmse} is not reduced relative to
E0. The aggregate gain is therefore driven mainly by improvements in
other portions of the hydrograph, whereas dry-season simulation remains
structurally difficult. This is consistent with recession theory: a
single linear reservoir cannot reproduce the effects of heterogeneous
hillslope storage, wetlands, aquifer connectivity and spatially variable
release times \citep{harman2009power,van2013influence}. It also helps
explain why the most substantial low-flow improvement appears only after
hybrid correction with ensemble quantiles. In other words, the baseflow
module adds physically meaningful memory, but it does not by itself
solve the low-flow problem.

This result has implications for calibration. A global NSE objective is
known to prioritize high-flow variance and can hide errors in the lower
part of the flow-duration curve \citep{pushpalatha2012review}. Future
versions should therefore include explicit low-flow criteria, such as
inverse-flow NSE or regime-weighted losses, rather than relying only on
global efficiency metrics. The broader ML literature on rainfall--runoff
prediction also supports the importance of antecedent runoff or memory
features for improving short-horizon runoff estimates
\citep{adnan2021comparison,okkan2021embedding}. For applications where
antecedent observed discharge is unavailable, physically simulated
memory descriptors may offer a practical alternative, but their value
must be tested against low-flow-specific metrics.

\subsection{L\'{e}vy perturbations: useful ensemble diversity, not calibrated uncertainty}
\label{subsec:discussion-levy}

The L\'{e}vy component should be interpreted cautiously. Heavy-tailed
perturbations are conceptually attractive in hydroclimatic settings
where rainfall intermittency, extreme events and nonlinear runoff
responses generate asymmetric response distributions. However, the
current results show that adding L\'{e}vy perturbations to the physical
model does not improve deterministic point-prediction skill. The
E3--E1 contrast slightly reduces NSE and increases RMSE, despite a
small gain in KGE. The perturbations therefore do not function as a
standalone point-prediction improvement mechanism in the current
configuration.

The uncertainty diagnostics are more restrictive. The nominal physical
bands are strongly underdispersive, with PICP far below the nominal
0.95 target. This means that the ensemble spread is too narrow to be
used as calibrated predictive uncertainty. In probabilistic hydrology,
sharp intervals are valuable only when they are also reliable; scoring
or reporting narrow bands without adequate coverage can be misleading
\citep{atger1999skill,franz2011evaluating,papacharalampous2022review}. Similar underdispersion has been reported in hydrological ensemble
prediction when forcing, parameter, structural and initial-state
uncertainties are not represented jointly
\citep{demirel2013effect,thiboult2016accounting,crochemore2016bias}.
StoHyMoLAP currently perturbs the physical response but does not yet
provide a full uncertainty propagation and post-processing chain.

The negative uncertainty result does not make the stochastic ensemble
irrelevant. It changes the role assigned to it. The ensemble is not yet
a reliable uncertainty product, but its quantiles remain informative
features for deterministic or conditional prediction. This distinction
is central to the paper: stochastic simulation can fail as a calibrated
interval while still creating useful descriptors of physical-response
asymmetry, spread and tail behaviour. The hydrological value of the
L\'{e}vy component is therefore mediated by the ML hybrid layer, not by
direct interval coverage.

\subsection{Why ensemble quantiles improve the hybrid model}
\label{subsec:discussion-hybrid}

The advantage of quantiles over ensemble means is the strongest and
most defensible empirical result. An ensemble mean compresses all
Monte Carlo members into a single central trajectory. This may be
adequate when the conditional response distribution is approximately
symmetric and homoscedastic, but it discards information about state-
dependent spread, skewness and tail behaviour. In the Ramis basin,
these properties are expected to vary between wet-season peaks and
dry-season recession. Quantile features retain this variation by
representing several points of the conditional distribution rather than
only its first moment.

This mechanism also explains why the quantile hybrid outperforms the
pure forcing-based ML baseline. Observed forcings describe atmospheric
inputs, but they do not directly encode the catchment response state
produced by the rainfall--runoff transformation. The physical ensemble
acts as a hydrological feature generator: it converts meteorological
forcing into discharge-consistent descriptors before ML correction.
Similar physical-output enrichment strategies have been reported in
hybrid hydrological modelling, where process-model outputs or synthetic
physical simulations provide state information that is difficult for a
data-driven model to infer from raw forcings alone
\citep{mekonnen2015hybrid,yang2020physical,xie2021physics,li2024process,zhu2025hybrid}. The present results extend that idea by showing that
quantiles of the stochastic physical response are more informative than
the ensemble mean alone.

Nevertheless, the interpretation should remain conservative. The
E6--E4 comparison combines two differences: the use of physical-ensemble
features and the use of distributional descriptors. It does not by
itself prove that all of the gain comes from physical structure rather
than from temporal aggregation or feature engineering. A useful future
ablation would compare ensemble quantiles against purely statistical lag
quantiles derived from observed forcings or from simple persistence
features. Such an experiment would separate the value of physical
simulation from the value of quantile summarization.

\subsection{Backend-aware interpretation of the sequential experiments}
\label{subsec:discussion-backend}

The numerical ranking places E8 first, but the backend audit prevents a
standard recurrent-network interpretation. E7 and E8 were labelled as
GRU experiments, yet the executed backend was a fallback multilayer
perceptron applied to flattened sequences. These experiments therefore
do not test the gating and hidden-state dynamics that define recurrent
architectures. They should be described as GRU-labelled sequential
hybrids under fallback-MLP execution, not as evidence that GRUs or other
recurrent neural networks outperform gradient boosting
\citep{kratzert2018rainfall,kratzert2019towards}.

The audit-adjusted interpretation is still useful. Quantile descriptors
outperform mean descriptors in both executed emulator families: E6
exceeds E5 in the gradient-boosting family, and E8 exceeds E7 in the
fallback sequential-feature family. This repeated ordering suggests
that descriptor choice is more important than the nominal architecture
label in the present evidence package. However, architecture-agnostic
performance should be treated as a hypothesis, not as a final claim. A
publication-ready extension should rerun E7 and E8 with an actual
TensorFlow or PyTorch recurrent backend and compare the resulting
architecture effect with the descriptor effect already measured here.

\subsection{Reproducibility, leakage control and submission readiness}
\label{subsec:discussion-leakage}

The audit tables strengthen the evidential value of the ablation by
making the modelling workflow inspectable. The current evidence package
reports strict temporal splitting, non-negative lags, sequence
construction within temporal subsets, train-only scalers and train-only
regime thresholds. These controls address common sources of optimistic
performance in hydrological ML, especially when temporal dependence and
lagged predictors are present. The audit therefore supports the main
comparative finding: the improvement of E6 over E5 and E4 is unlikely
to be an artefact of obvious temporal leakage.

At the same time, the WARN status should remain visible. WARN does not
mean that leakage was detected, but it does identify issues that affect
scientific interpretation: missing observed-flow values, low-flow
physical limitations, underdispersive stochastic intervals, optional
figure entries and fallback ML backends. Keeping these warnings in the
main manuscript or supplementary material is preferable to hiding them,
because it clarifies which conclusions are robust and which require
additional computation. This position is aligned with calls for more
transparent, reproducible and auditable hydrological modelling
\citep{hall2021hydrologist,lees2021benchmarking,arsenault2022continuous}.

For submission to \emph{Journal of Hydrology}, the strongest claim is
therefore deliberately narrow: quantile descriptors extracted from a
stochastic physical ensemble improve backend-consistent hybrid
simulation of daily discharge in the Ramis basin. Claims about
calibrated uncertainty, recurrent-network superiority or operational
transferability should remain conditional until the uncertainty bands
are recalibrated, the true recurrent backend is rerun and the final data
provenance is fully documented.


